%% 首屏主结果图：一份报告的构建耗时。
%% 数据取自本机 TeX Live 2026 上 latexmk 的两次完整构建（含 bibtex），
%% 换机器后请重测并更新坐标，不要沿用旧数字。
\begin{figure}[htbp]
  \centering
  \begin{tikzpicture}
    \begin{axis}[
      ybar,
      bar width=26pt,
      width=0.86\textwidth,
      height=5.6cm,
      enlarge x limits=0.28,
      ymin=0,
      ymax=20,
      ylabel={构建耗时（秒）},
      ylabel style={font=\small},
      symbolic x coords={英文示例, 中文示例, 空白骨架},
      xtick=data,
      xticklabel style={font=\small},
      tick style={draw=none},
      ymajorgrids,
      grid style={bgrule!45},
      axis line style={bgrule},
      nodes near coords,
      nodes near coords style={font=\scriptsize, color=bgink},
    ]
      \addplot[fill=bgbrand, draw=none] coordinates
        {(英文示例,8.5) (中文示例,15.7) (空白骨架,1.3)};
    \end{axis}
  \end{tikzpicture}
  \caption{本机完整构建耗时（两遍编译 + 参考文献）。中文示例因为要走 XeLaTeX
  与 ctex，约为英文示例的两倍，但仍然在一次编译停顿之内。}
  \label{fig:hero}
\end{figure}
